\subsection{Discrete-Input Mutual Information}
\label{subsec:discrete_input_mi}

Let $X\in\mathcal{X}=\{\alpha_1,\ldots,\alpha_M\}$ denote the transmitted
symbol, where $M=256$, $\Pr[X=\alpha_m]=p_m$, and $\sum_{m=1}^{M}p_m=1$.
The normalized constellation satisfies
\begin{equation}
    \sum_{m=1}^{M}p_m\alpha_m=0,
    \qquad
    2\sum_{m=1}^{M}p_m|\alpha_m|^2=V_A.
\end{equation}
For an instantaneous channel transmittance $T$, the complex baseband auxiliary
channel is
\begin{equation}
    Y=\sqrt{T}\alpha_X+N,
    \qquad
    N\sim\mathcal{CN}(0,\sigma_c^2),
    \qquad
    \sigma_c^2=1+\frac{T\epsilon}{2}.
\end{equation}
Thus, $\Re\{N\}$ and $\Im\{N\}$ are independent zero-mean Gaussian random
variables with variance $\sigma_c^2/2$. This convention gives
\begin{equation}
    \mathsf{SNR}
    =\frac{T V_A/2}{1+T\epsilon/2}
    =\frac{T V_A}{2+T\epsilon},
\end{equation}
which is identical to the signal-to-noise convention used by the Gaussian-input
benchmark.

The discrete-input mutual information is
\begin{equation}
    I_{AB}=I(X;Y)=H(X)-H(X|Y).
\end{equation}
For an auxiliary channel $Q_{Y|X}$, an achievable mismatched-decoding rate is
\begin{equation}
    I_{AB}\geq R_{\mathrm{sym}}
    =H(X)+\mathbb{E}_{X,Y}
    \left[\log_2 Q_{X|Y}(X|Y)\right].
\end{equation}
Using the unnormalized likelihood
$Q_{Y|X}(y|\alpha_m)\propto
\exp[-|y-\sqrt{T}\alpha_m|^2/\sigma_c^2]$, the auxiliary posterior is
\begin{equation}
    Q_{X|Y}(\alpha_m|y)
    =\frac{p_m\exp\!\left(-|y-\sqrt{T}\alpha_m|^2/\sigma_c^2\right)}
    {\sum_{j=1}^{M}p_j\exp\!\left(-|y-\sqrt{T}\alpha_j|^2/\sigma_c^2\right)}.
\end{equation}
The likelihood normalization is omitted because it cancels in the posterior.
The posterior is evaluated numerically in the log domain using the
log-sum-exp operation.

For samples $(x_i,y_i)_{i=1}^{N}$, the Monte Carlo estimator is
\begin{align}
    \widehat{R}_{\mathrm{sym}}
    ={}&-\sum_{m=1}^{M}p_m\log_2 p_m \nonumber\\
    &+\frac{1}{N}\sum_{i=1}^{N}\log_2\!\left[
    \frac{p_{x_i}\exp\!\left(-|y_i-\sqrt{T}\alpha_{x_i}|^2/\sigma_c^2\right)}
    {\sum_{j=1}^{M}p_j\exp\!\left(-|y_i-\sqrt{T}\alpha_j|^2/\sigma_c^2\right)}
    \right].
\end{align}
In the implementation, all transmitted constellation points are enumerated
exactly and weighted by $p_m$, while only the complex additive noise is sampled.
For fading samples $T_b$, the estimator is evaluated conditionally for every
$T_b$ and subsequently averaged; the mean transmittance is not substituted into
the nonlinear mutual-information expression.

The raw asymptotic secret key rate is therefore evaluated as
\begin{equation}
    K_{\mathrm{raw}}=\beta\widehat{R}_{\mathrm{sym}}-\chi_{BE}.
\end{equation}
No nonnegativity projection is applied during optimization. Only reported
curves and tables use $K_{\mathrm{reported}}=\max(K_{\mathrm{raw}},0)$.
The Holevo-information calculation $\chi_{BE}$ and its covariance-matrix model
remain unchanged.

For comparison only, the Gaussian-input reference mutual information is
\begin{equation}
    I_{AB}^{\mathrm{G}}
    =\log_2\!\left(1+\frac{T V_A}{2+T\epsilon}\right).
\end{equation}
The quantity $I_{AB}^{\mathrm{G}}$ is not used to optimize the probabilistic
shaping network. Unlike the Gaussian reference, the finite-alphabet achievable
rate saturates at $H(X)\leq\log_2(256)=8$ bits per symbol.
